% Figure: blowdown wind-tunnel reservoir pressure decay. A fixed-volume
% high-pressure tank feeds a choked test section; the reservoir pressure
% falls exponentially as mass leaves at the choked rate, and the usable
% run ends when the pressure can no longer sustain the design test-section
% Mach number (the minimum-pressure floor). Schematic with a representative
% decay curve and the run-time window marked.
\begin{figure}[htbp]
\centering
\begin{tikzpicture}
\begin{axis}[
  width=12cm, height=8cm,
  xlabel={time $t$ (s)},
  ylabel={reservoir pressure $p_0(t)$ (bar)},
  xmin=0, xmax=60, ymin=0, ymax=22,
  axis lines=left,
  grid=major, grid style={enggray!20},
  tick label style={font=\scriptsize},
  label style={font=\small},
  legend style={font=\scriptsize, at={(0.97,0.97)}, anchor=north east},
]
% exponential decay p0(t) = p0_init * exp(-t/tau), tau = 25 s
\addplot[engblue, very thick, domain=0:60, samples=200] {
  20*exp(-x/25) };
\addlegendentry{$p_0(t)=p_{0,i}\,e^{-t/\tau}$}
% minimum-pressure floor to hold the design Mach number: 8 bar
\addplot[engred, dashed, domain=0:60, samples=2] {8};
\node[engred, font=\scriptsize, anchor=south west] at (axis cs:0.5,8.1) {minimum $p_0$ to hold design $M$};
% run ends where curve crosses the floor: 20 exp(-t/25)=8 -> t = 25 ln(2.5) = 22.9 s
\addplot[engblack, only marks, mark=*, mark options={fill=engblack}] coordinates {(22.9,8)};
\node[engblack, font=\scriptsize, anchor=south west] at (axis cs:23.5,8.6) {run ends, $t\approx 23$ s};
% shade the usable run window
\draw[englime, very thick] (axis cs:0,1.0) -- (axis cs:22.9,1.0);
\node[englime, font=\scriptsize, anchor=south] at (axis cs:11.5,1.1) {usable run};
\end{axis}
\end{tikzpicture}
\caption[Blowdown reservoir pressure decay]{Reservoir pressure of a
fixed-volume blowdown wind tunnel feeding a choked test section. With
the throat choked, mass leaves at a rate proportional to the
instantaneous reservoir pressure, so for an isothermal reservoir the
pressure decays exponentially, $p_0(t)=p_{0,i}\,e^{-t/\tau}$, with the
time constant $\tau$ set by the tank volume divided by the choked
volumetric efflux. The run ends not when the tank empties but when the
reservoir pressure falls below the floor needed to keep the test-section
pressure ratio above its design value (here $8$ bar from an initial
$20$ bar), which fixes the usable run time at $t=\tau\ln(p_{0,i}/p_{0,
\min})\approx 23\,$s for the curve shown. The blowdown vignette in the
text carries this arithmetic for a sized tank.}
\label{fig:vol03ch13:blowdown}
\end{figure}
